\documentclass{article}

\usepackage{amsfonts}
\usepackage{amsmath}

\title{Controlli Automatici}
\author{Leonardo Florenzano}

\begin{document}

\maketitle

\section{Modelli e Algoritmi}

L'obbiettivo del corso è di sostituire il ragionamento umano, perciò l'elemento che controlla.
Possiamo rappresentare la realtà come un modello, più o meno preciso.
Una rete di controllo ha:

\begin{itemize}
\item sensori per avere una immagine dello stato del modello.
\item una uscita per pilotare il modello in base alla logica di controllo.
\end{itemize}

Lo stato interno del modello si può rappresentare con strumenti matematici, come le equazioni differenziali ordinarie:

\[ \dot{x}(t) = f(x(t), u(t), t) \]

L'uscista della rete invece si può rappresentare con una semplice trasformazione:

\[ y = h(x(t), t) \]

Diamo per vero che entrambe le funzioni di stato e uscita sono lineari.
Si può generalizzare ancora di più la descrizione del modello, espandendo le funzioni \(u(t), x(t), y(t)\) a più dimensioni:

\begin{align*}
  u(t) &= \begin{bmatrix} u_1(t) \\ u_2(t) \\ \dots \\ u_n(t) \end{bmatrix} \in \mathbb{R}^n\\
  x(t) &= \begin{bmatrix} x_1(t) \\ x_2(t) \\ \dots \\ x_n(t) \end{bmatrix} \in \mathbb{R}^m\\
  y(t) &= \begin{bmatrix} y_1(t) \\ y_2(t) \\ \dots \\ y_n(t) \end{bmatrix} \in \mathbb{R}^o
\end{align*}

Possiamo definire \(\dot{x}(t) = f(x(t), u(t), t)\) come \textbf{traiettoria del sistema}.
In particolare \(x(t)\) si dice traiettoria dello stato e \(y(t)\) si dice traiettoria dell'uscita.
Sistemi che non hanno un ingresso, vengono detti \textbf{non forzati}.

\vspace{10pt}

Un sistema lineare tempo invariante si può scrivere nella seguente maniera:

\begin{align*}
  \dot{x}(t) &= A(t)x(t) + B(t)u(t)\\
  y(t) &= C(t)x(t) + D(t)u(t)
\end{align*}

Dove \(A(t), B(t), C(t), D(t)\) sono delle matrici.

\[
  \begin{bmatrix} \dot{x}_1(t) \\ \vdots \\ \dot{x}_n(t) \end{bmatrix}
  =
  \begin{bmatrix}
    a_{11} & \dots & a_{1n}\\
    \vdots & \ddots & \vdots\\
    a_{m1} & \dots & a_{mn}
  \end{bmatrix} 
  \begin{bmatrix} x_1(t) \\ \vdots \\ x_n(t) \end{bmatrix}
  +
  \begin{bmatrix}
    b_{11} & \dots & b_{1n}\\
    \vdots & \ddots & \vdots\\
    b_{l1} & \dots & b_{ln}
  \end{bmatrix}
  \begin{bmatrix} u_1(t) \\ \vdots \\ u_n(t) \end{bmatrix}
\]

La traiettoria di stato \(x(t)\) può essere come somma di due componenti

\[ x(t) = x_l(t) + x_f(t) \]

Queste due componenti sono chiamate

\begin{itemize}
\item Evoluzione libera: \(x_l(t) = A x(t) , x(0) = x_0 \neq 0\)
\item Evoluzione forzata: \(x_f(t) = B u(t) , u(0) = u_0 \neq 0\)
\end{itemize}

Se io prendo uno stato \(x_E=0\) e un ingresso \(u_E=0\), porto il sistema in equilibrio.

\end{document}
